\documentclass[10pt,a4paper]{amsart}
\usepackage{amsmath,amssymb,amsthm,graphicx}
\usepackage{verbatim}
\newenvironment{reference}{\comment}{\endcomment}
%\newenvironment{reference}{}{}
\newenvironment{slogan}{\comment}{\endcomment}
\newenvironment{history}{\comment}{\endcomment}
% For commutative diagrams we use Xy-pic
\usepackage[all]{xy}

% We use 2cell for 2-commutative diagrams.
\xyoption{2cell}
% Theorem environments.
%
\theoremstyle{plain}
\newtheorem{theorem}[subsection]{Theorem}
\newtheorem{proposition}[subsection]{Proposition}
\newtheorem{lemma}[subsection]{Lemma}

\theoremstyle{definition}
\newtheorem{definition}[subsection]{Definition}
\newtheorem{example}[subsection]{Example}
\newtheorem{exercise}[subsection]{Exercise}
\newtheorem{situation}[subsection]{Situation}

\theoremstyle{remark}
\newtheorem{remark}[subsection]{Remark}
\newtheorem{remarks}[subsection]{Remarks}

\numberwithin{equation}{subsection}

% Macros
%
\def\lim{\mathop{\mathrm{lim}}\nolimits}
\def\colim{\mathop{\mathrm{colim}}\nolimits}
\def\Spec{\mathop{\mathrm{Spec}}}
\def\Hom{\mathop{\mathrm{Hom}}\nolimits}
\def\Ext{\mathop{\mathrm{Ext}}\nolimits}
\def\SheafHom{\mathop{\mathcal{H}\!\mathit{om}}\nolimits}
\def\SheafExt{\mathop{\mathcal{E}\!\mathit{xt}}\nolimits}
\def\Sch{\mathit{Sch}}
\def\Mor{\mathop{\mathrm{Mor}}\nolimits}
\def\Ob{\mathop{\mathrm{Ob}}\nolimits}
\def\Sh{\mathop{\mathit{Sh}}\nolimits}
\def\NL{\mathop{N\!L}\nolimits}
\def\CH{\mathop{\mathrm{CH}}\nolimits}
\def\proetale{{pro\text{-}\acute{e}tale}}
\def\etale{{\acute{e}tale}}
\def\QCoh{\mathit{QCoh}}
\def\Ker{\mathop{\mathrm{Ker}}}
\def\Im{\mathop{\mathrm{Im}}}
\def\Coker{\mathop{\mathrm{Coker}}}
\def\Coim{\mathop{\mathrm{Coim}}}

% Boxtimes
%
\DeclareMathSymbol{\boxtimes}{\mathbin}{AMSa}{"02}

%
% Macros for moduli stacks/spaces
%
\def\QCohstack{\mathcal{QC}\!\mathit{oh}}
\def\Cohstack{\mathcal{C}\!\mathit{oh}}
\def\Spacesstack{\mathcal{S}\!\mathit{paces}}
\def\Quotfunctor{\mathrm{Quot}}
\def\Hilbfunctor{\mathrm{Hilb}}
\def\Curvesstack{\mathcal{C}\!\mathit{urves}}
\def\Polarizedstack{\mathcal{P}\!\mathit{olarized}}
\def\Complexesstack{\mathcal{C}\!\mathit{omplexes}}
% \Pic is the operator that assigns to X its picard group, usage \Pic(X)
% \Picardstack_{X/B} denotes the Picard stack of X over B
% \Picardfunctor_{X/B} denotes the Picard functor of X over B
\def\Pic{\mathop{\mathrm{Pic}}\nolimits}
\def\Picardstack{\mathcal{P}\!\mathit{ic}}
\def\Picardfunctor{\mathrm{Pic}}
\def\Deformationcategory{\mathcal{D}\!\mathit{ef}}
\usepackage{cleveref}
\crefname{originalreference}{original reference}{original references}
\Crefname{originalreference}{Original reference}{Original references}
\makeatletter
\newlabel{algebra-lemma-find-sequence-image-regular}{{algebra-lemma-find-sequence-image-regular}{source}{}{sourcepointer-1}{}}
\newlabel{algebra-lemma-find-sequence-image-regular@cref}{{[originalreference][0][]algebra-lemma-find-sequence-image-regular}{[1][0][]source}{}{}{}}
\makeatother
\crefname{originalreference}{original reference}{original references}
\Crefname{originalreference}{Original reference}{Original references}
\makeatletter
\newlabel{algebra-lemma-regular-ring-CM}{{algebra-lemma-regular-ring-CM}{source}{}{sourcepointer-2}{}}
\newlabel{algebra-lemma-regular-ring-CM@cref}{{[originalreference][0][]algebra-lemma-regular-ring-CM}{[1][0][]source}{}{}{}}
\makeatother
\crefname{originallemma}{original lemma}{original lemmas}
\Crefname{originallemma}{Original lemma}{Original lemmas}
\makeatletter
\newlabel{lemma-enough-finite-type-points}{{lemma-enough-finite-type-points}{source}{}{sourcepointer-3}{}}
\newlabel{lemma-enough-finite-type-points@cref}{{[originallemma][0][]lemma-enough-finite-type-points}{[1][0][]source}{}{}{}}
\makeatother
\crefname{originallemma}{original lemma}{original lemmas}
\Crefname{originallemma}{Original lemma}{Original lemmas}
\makeatletter
\newlabel{lemma-point-finite-type}{{lemma-point-finite-type}{source}{}{sourcepointer-4}{}}
\newlabel{lemma-point-finite-type@cref}{{[originallemma][0][]lemma-point-finite-type}{[1][0][]source}{}{}{}}
\makeatother
\crefname{originallemma}{original lemma}{original lemmas}
\Crefname{originallemma}{Original lemma}{Original lemmas}
\makeatletter
\newlabel{lemma-point-finite-type-monomorphism}{{lemma-point-finite-type-monomorphism}{source}{}{sourcepointer-5}{}}
\newlabel{lemma-point-finite-type-monomorphism@cref}{{[originallemma][0][]lemma-point-finite-type-monomorphism}{[1][0][]source}{}{}{}}
\makeatother
\crefname{originallemma}{original lemma}{original lemmas}
\Crefname{originallemma}{Original lemma}{Original lemmas}
\makeatletter
\newlabel{lemma-properties-covering-imply-diagonal}{{lemma-properties-covering-imply-diagonal}{source}{}{sourcepointer-6}{}}
\newlabel{lemma-properties-covering-imply-diagonal@cref}{{[originallemma][0][]lemma-properties-covering-imply-diagonal}{[1][0][]source}{}{}{}}
\makeatother
\crefname{originallemma}{original lemma}{original lemmas}
\Crefname{originallemma}{Original lemma}{Original lemmas}
\makeatletter
\newlabel{lemma-quasi-finite-at-point}{{lemma-quasi-finite-at-point}{source}{}{sourcepointer-7}{}}
\newlabel{lemma-quasi-finite-at-point@cref}{{[originallemma][0][]lemma-quasi-finite-at-point}{[1][0][]source}{}{}{}}
\makeatother
\crefname{originallemma}{original lemma}{original lemmas}
\Crefname{originallemma}{Original lemma}{Original lemmas}
\makeatletter
\newlabel{lemma-slice}{{lemma-slice}{source}{}{sourcepointer-8}{}}
\newlabel{lemma-slice@cref}{{[originallemma][0][]lemma-slice}{[1][0][]source}{}{}{}}
\makeatother
\crefname{originalreference}{original reference}{original references}
\Crefname{originalreference}{Original reference}{Original references}
\makeatletter
\newlabel{spaces-groupoids-definition-stabilizer-groupoid}{{spaces-groupoids-definition-stabilizer-groupoid}{source}{}{sourcepointer-9}{}}
\newlabel{spaces-groupoids-definition-stabilizer-groupoid@cref}{{[originalreference][0][]spaces-groupoids-definition-stabilizer-groupoid}{[1][0][]source}{}{}{}}
\makeatother
\crefname{originalreference}{original reference}{original references}
\Crefname{originalreference}{Original reference}{Original references}
\makeatletter
\newlabel{spaces-more-groupoids-lemma-groupoid-on-field-dimension-equal-stabilizer}{{spaces-more-groupoids-lemma-groupoid-on-field-dimension-equal-stabilizer}{source}{}{sourcepointer-10}{}}
\newlabel{spaces-more-groupoids-lemma-groupoid-on-field-dimension-equal-stabilizer@cref}{{[originalreference][0][]spaces-more-groupoids-lemma-groupoid-on-field-dimension-equal-stabilizer}{[1][0][]source}{}{}{}}
\makeatother
\crefname{originalreference}{original reference}{original references}
\Crefname{originalreference}{Original reference}{Original references}
\makeatletter
\newlabel{spaces-morphisms-lemma-dimension-fibre-at-a-point}{{spaces-morphisms-lemma-dimension-fibre-at-a-point}{source}{}{sourcepointer-11}{}}
\newlabel{spaces-morphisms-lemma-dimension-fibre-at-a-point@cref}{{[originalreference][0][]spaces-morphisms-lemma-dimension-fibre-at-a-point}{[1][0][]source}{}{}{}}
\makeatother
\crefname{originalreference}{original reference}{original references}
\Crefname{originalreference}{Original reference}{Original references}
\makeatletter
\newlabel{spaces-over-fields-lemma-smooth-regular}{{spaces-over-fields-lemma-smooth-regular}{source}{}{sourcepointer-12}{}}
\newlabel{spaces-over-fields-lemma-smooth-regular@cref}{{[originalreference][0][]spaces-over-fields-lemma-smooth-regular}{[1][0][]source}{}{}{}}
\makeatother
\crefname{originalreference}{original reference}{original references}
\Crefname{originalreference}{Original reference}{Original references}
\makeatletter
\newlabel{spaces-properties-lemma-regular}{{spaces-properties-lemma-regular}{source}{}{sourcepointer-13}{}}
\newlabel{spaces-properties-lemma-regular@cref}{{[originalreference][0][]spaces-properties-lemma-regular}{[1][0][]source}{}{}{}}
\makeatother
\crefname{originalreference}{original reference}{original references}
\Crefname{originalreference}{Original reference}{Original references}
\makeatletter
\newlabel{stacks-properties-lemma-characterize-surjective}{{stacks-properties-lemma-characterize-surjective}{source}{}{sourcepointer-14}{}}
\newlabel{stacks-properties-lemma-characterize-surjective@cref}{{[originalreference][0][]stacks-properties-lemma-characterize-surjective}{[1][0][]source}{}{}{}}
\makeatother
\crefname{originalreference}{original reference}{original references}
\Crefname{originalreference}{Original reference}{Original references}
\makeatletter
\newlabel{stacks-properties-lemma-topology-points}{{stacks-properties-lemma-topology-points}{source}{}{sourcepointer-15}{}}
\newlabel{stacks-properties-lemma-topology-points@cref}{{[originalreference][0][]stacks-properties-lemma-topology-points}{[1][0][]source}{}{}{}}
\makeatother
\title{A quasi-DM stack is exactly an algebraic stack admitting a locally quasi-finite flat cover by a scheme}
\author{The Stacks Project Authors}
\date{stacks-project:stacks-morphisms.tex:local-snapshot-20261012; GNU FDL 1.2 or later}
\begin{document}
\maketitle
\noindent\textit{Editorial note.} All delivered mathematical text is the author's own, copied from the pinned Stacks Project chapter source sources/stacks-algebraic-stacks-morphisms/pinned-original-source.tex; no statement or proof step was rewritten, reordered or removed.\par\medskip
\noindent\textit{Editorial note.} Preamble provenance: the Stacks Project keeps one single shared preamble for every chapter; this card takes it directly from the staged sources/stacks-algebraic-stacks-morphisms/preamble.tex as three verbatim spans (lines 8-12, 14-18 and 154-217), so no derived file was written and nothing was added or redefined. The wrapper class and the base packages are supplied by the assembler.\par\medskip
\noindent\textit{Editorial note.} Bibliography provenance: neither the counted statement nor its proof contains a citation (this lane rescans the whole span before writing the card), so the card carries no bibliography and no entry was reconstructed.\par\medskip
\noindent\textit{Editorial note.} Selection provenance: this card was opened under the file-based control contract (CONTROL.md), with the mechanical pre-card gate run on its statement and proof interval before the card was written.\par\medskip
\noindent\textit{Allowed standard prerequisites (editorial).}\begin{itemize}
\item The residual gerbe at a finite type point, together with the description of quasi-DM points through the automorphism group being locally quasi-finite. Invoked inside the counted proof (source lines 4296-4310). Reduces the problem to the one-point case where the automorphism group of the point is locally quasi-finite.
\item The structure of a one-point stack through a group scheme acting on a scheme, with the point corresponding to a section. Used inside the counted proof (source lines 4311-4340). Presents the residual gerbe by a group action so that the automorphism group becomes a group algebraic space over the residue field.
\item The statement that a locally quasi-finite group algebraic space over a field is a scheme with finite type structure, and the localization of a locally quasi-finite morphism at a point. Invoked inside the counted proof (source lines 4341-4380) as cross-chapter pointers. Produces a finite flat group scheme map through which the cover is built.
\item The criterion for local quasi-finiteness through discrete field fibres. Invoked near the end of the counted proof (source lines 4381-4400). Verifies that the constructed cover is locally quasi-finite.
\end{itemize}\medskip
\ifdefined\hypertarget\hypertarget{sourcepointer-1}{}\fi
\noindent\textit{Source reference.} Original reference algebra-lemma-find-sequence-image-regular: sources/stacks-algebraic-stacks-morphisms/pinned-original-source.tex, source lines 4367--4367. Navigation pointer only: the counted text refers to the Stacks item algebra-lemma-find-sequence-image-regular, which belongs to another chapter of the project and is not part of this card. The pointer carries the target's own label as its number because this chapter's staged snapshot does not contain the other chapter's tag table; its content is not reproduced here.\par\medskip
\ifdefined\hypertarget\hypertarget{sourcepointer-2}{}\fi
\noindent\textit{Source reference.} Original reference algebra-lemma-regular-ring-CM: sources/stacks-algebraic-stacks-morphisms/pinned-original-source.tex, source lines 4365--4365. Navigation pointer only: the counted text refers to the Stacks item algebra-lemma-regular-ring-CM, which belongs to another chapter of the project and is not part of this card. The pointer carries the target's own label as its number because this chapter's staged snapshot does not contain the other chapter's tag table; its content is not reproduced here.\par\medskip
\ifdefined\hypertarget\hypertarget{sourcepointer-3}{}\fi
\noindent\textit{Source reference.} Original lemma lemma-enough-finite-type-points: sources/stacks-algebraic-stacks-morphisms/pinned-original-source.tex, source lines 3767--3767. Navigation pointer only: the counted text refers to the Stacks item lemma-enough-finite-type-points, which is an existing labeled result of this same chapter and is not part of this card. The pointer carries the target's own label as its number and its definition line as its span; its content is not reproduced here.\par\medskip
\ifdefined\hypertarget\hypertarget{sourcepointer-4}{}\fi
\noindent\textit{Source reference.} Original lemma lemma-point-finite-type: sources/stacks-algebraic-stacks-morphisms/pinned-original-source.tex, source lines 3656--3656. Navigation pointer only: the counted text refers to the Stacks item lemma-point-finite-type, which is an existing labeled result of this same chapter and is not part of this card. The pointer carries the target's own label as its number and its definition line as its span; its content is not reproduced here.\par\medskip
\ifdefined\hypertarget\hypertarget{sourcepointer-5}{}\fi
\noindent\textit{Source reference.} Original lemma lemma-point-finite-type-monomorphism: sources/stacks-algebraic-stacks-morphisms/pinned-original-source.tex, source lines 3803--3803. Navigation pointer only: the counted text refers to the Stacks item lemma-point-finite-type-monomorphism, which is an existing labeled result of this same chapter and is not part of this card. The pointer carries the target's own label as its number and its definition line as its span; its content is not reproduced here.\par\medskip
\ifdefined\hypertarget\hypertarget{sourcepointer-6}{}\fi
\noindent\textit{Source reference.} Original lemma lemma-properties-covering-imply-diagonal: sources/stacks-algebraic-stacks-morphisms/pinned-original-source.tex, source lines 1117--1117. Navigation pointer only: the counted text refers to the Stacks item lemma-properties-covering-imply-diagonal, which is an existing labeled result of this same chapter and is not part of this card. The pointer carries the target's own label as its number and its definition line as its span; its content is not reproduced here.\par\medskip
\ifdefined\hypertarget\hypertarget{sourcepointer-7}{}\fi
\noindent\textit{Source reference.} Original lemma lemma-quasi-finite-at-point: sources/stacks-algebraic-stacks-morphisms/pinned-original-source.tex, source lines 4248--4248. Navigation pointer only: the counted text refers to the Stacks item lemma-quasi-finite-at-point, which is an existing labeled result of this same chapter and is not part of this card. The pointer carries the target's own label as its number and its definition line as its span; its content is not reproduced here.\par\medskip
\ifdefined\hypertarget\hypertarget{sourcepointer-8}{}\fi
\noindent\textit{Source reference.} Original lemma lemma-slice: sources/stacks-algebraic-stacks-morphisms/pinned-original-source.tex, source lines 4177--4177. Navigation pointer only: the counted text refers to the Stacks item lemma-slice, which is an existing labeled result of this same chapter and is not part of this card. The pointer carries the target's own label as its number and its definition line as its span; its content is not reproduced here.\par\medskip
\ifdefined\hypertarget\hypertarget{sourcepointer-9}{}\fi
\noindent\textit{Source reference.} Original reference spaces-groupoids-definition-stabilizer-groupoid: sources/stacks-algebraic-stacks-morphisms/pinned-original-source.tex, source lines 4323--4323. Navigation pointer only: the counted text refers to the Stacks item spaces-groupoids-definition-stabilizer-groupoid, which belongs to another chapter of the project and is not part of this card. The pointer carries the target's own label as its number because this chapter's staged snapshot does not contain the other chapter's tag table; its content is not reproduced here.\par\medskip
\ifdefined\hypertarget\hypertarget{sourcepointer-10}{}\fi
\noindent\textit{Source reference.} Original reference spaces-more-groupoids-lemma-groupoid-on-field-dimension-equal-stabilizer: sources/stacks-algebraic-stacks-morphisms/pinned-original-source.tex, source lines 4333--4333. Navigation pointer only: the counted text refers to the Stacks item spaces-more-groupoids-lemma-groupoid-on-field-dimension-equal-stabilizer, which belongs to another chapter of the project and is not part of this card. The pointer carries the target's own label as its number because this chapter's staged snapshot does not contain the other chapter's tag table; its content is not reproduced here.\par\medskip
\ifdefined\hypertarget\hypertarget{sourcepointer-11}{}\fi
\noindent\textit{Source reference.} Original reference spaces-morphisms-lemma-dimension-fibre-at-a-point: sources/stacks-algebraic-stacks-morphisms/pinned-original-source.tex, source lines 4383--4383. Navigation pointer only: the counted text refers to the Stacks item spaces-morphisms-lemma-dimension-fibre-at-a-point, which belongs to another chapter of the project and is not part of this card. The pointer carries the target's own label as its number because this chapter's staged snapshot does not contain the other chapter's tag table; its content is not reproduced here.\par\medskip
\ifdefined\hypertarget\hypertarget{sourcepointer-12}{}\fi
\noindent\textit{Source reference.} Original reference spaces-over-fields-lemma-smooth-regular: sources/stacks-algebraic-stacks-morphisms/pinned-original-source.tex, source lines 4361--4361. Navigation pointer only: the counted text refers to the Stacks item spaces-over-fields-lemma-smooth-regular, which belongs to another chapter of the project and is not part of this card. The pointer carries the target's own label as its number because this chapter's staged snapshot does not contain the other chapter's tag table; its content is not reproduced here.\par\medskip
\ifdefined\hypertarget\hypertarget{sourcepointer-13}{}\fi
\noindent\textit{Source reference.} Original reference spaces-properties-lemma-regular: sources/stacks-algebraic-stacks-morphisms/pinned-original-source.tex, source lines 4363--4363. Navigation pointer only: the counted text refers to the Stacks item spaces-properties-lemma-regular, which belongs to another chapter of the project and is not part of this card. The pointer carries the target's own label as its number because this chapter's staged snapshot does not contain the other chapter's tag table; its content is not reproduced here.\par\medskip
\ifdefined\hypertarget\hypertarget{sourcepointer-14}{}\fi
\noindent\textit{Source reference.} Original reference stacks-properties-lemma-characterize-surjective: sources/stacks-algebraic-stacks-morphisms/pinned-original-source.tex, source lines 4141--4141. Navigation pointer only: the counted text refers to the Stacks item stacks-properties-lemma-characterize-surjective, which belongs to another chapter of the project and is not part of this card. The pointer carries the target's own label as its number because this chapter's staged snapshot does not contain the other chapter's tag table; its content is not reproduced here.\par\medskip
\ifdefined\hypertarget\hypertarget{sourcepointer-15}{}\fi
\noindent\textit{Source reference.} Original reference stacks-properties-lemma-topology-points: sources/stacks-algebraic-stacks-morphisms/pinned-original-source.tex, source lines 2128--2128. Navigation pointer only: the counted text refers to the Stacks item stacks-properties-lemma-topology-points, which belongs to another chapter of the project and is not part of this card. The pointer carries the target's own label as its number because this chapter's staged snapshot does not contain the other chapter's tag table; its content is not reproduced here.\par\medskip
\begin{theorem}
\label{theorem-quasi-DM}
Let $\mathcal{X}$ be an algebraic stack. The following are equivalent
\begin{enumerate}
\item $\mathcal{X}$ is quasi-DM, and
\item there exists a scheme $W$ and a surjective, flat, locally finitely
presented, locally quasi-finite morphism $W \to \mathcal{X}$.
\end{enumerate}
\end{theorem}

\begin{proof}
The implication (2) $\Rightarrow$ (1) is
Lemma \ref{lemma-properties-covering-imply-diagonal}.
Assume (1).
Let $x \in |\mathcal{X}|$ be a finite type point. We will produce a scheme
over $\mathcal{X}$ which ``works'' in a neighbourhood of $x$. At the end
of the proof we will take the disjoint union of all of these to conclude.

\medskip\noindent
Let $U$ be an affine scheme, $U \to \mathcal{X}$ a smooth morphism, and
$u \in U$ a closed point which maps to $x$, see
Lemma \ref{lemma-point-finite-type}.
Denote $u = \Spec(\kappa(u))$ as usual. Consider the following
commutative diagram
$$
\xymatrix{
u \ar[d] & R \ar[l] \ar[d] \\
U \ar[d] & F \ar[d] \ar[l]^p \\
\mathcal{X} & u \ar[l]
}
$$
with both squares fibre product squares, in particular
$R = u \times_\mathcal{X} u$. In the proof of
Lemma \ref{lemma-point-finite-type-monomorphism}
we have seen that $(u, R, s, t, c)$ is a groupoid in algebraic spaces
with $s, t$ locally of finite type. Let $G \to u$ be the stabilizer group
algebraic space (see
Groupoids in Spaces, Definition
\ref{spaces-groupoids-definition-stabilizer-groupoid}).
Note that
$$
G = R \times_{(u \times u)} u =
(u \times_\mathcal{X} u) \times_{(u \times u)} u =
\mathcal{X} \times_{\mathcal{X} \times \mathcal{X}} u.
$$
As $\mathcal{X}$ is quasi-DM we see that
$G$ is locally quasi-finite over $u$. By
More on Groupoids in Spaces, Lemma
\ref{spaces-more-groupoids-lemma-groupoid-on-field-dimension-equal-stabilizer}
we have $\dim(R) = 0$.

\medskip\noindent
Let $e : u \to R$ be the identity of the groupoid. The composition
$u \to R \to u$ is equal to the identity morphism of $u$.
Note that $R \subset F$ is a closed
subspace as $u \subset U$ is a closed subscheme. Hence we can also think
of $e$ as a point of $F$; call it $z$.
Consider the maps of \'etale local rings
$$
\mathcal{O}_{U, u}
\xrightarrow{p^\sharp}
\mathcal{O}_{F, \overline{z}}
\longrightarrow
\mathcal{O}_{R, \overline{e}}
$$
Note that $\mathcal{O}_{R, \overline{e}}$ has dimension $0$ by the result
of the first paragraph. On the other hand, the kernel of the second arrow is
$p^\sharp(\mathfrak m_u)\mathcal{O}_{F, \overline{z}}$ as
$R$ is cut out in $F$ by $\mathfrak m_u$. Thus we see that
$$
\mathfrak m_{\overline{z}} =
\sqrt{p^\sharp(\mathfrak m_u)\mathcal{O}_{F, \overline{z}}}
$$
On the other hand, as the morphism $U \to \mathcal{X}$ is smooth
we see that $F \to u$ is a smooth morphism of algebraic spaces.
This means that $F$ is a regular algebraic space
(Spaces over Fields, Lemma \ref{spaces-over-fields-lemma-smooth-regular}).
Hence $\mathcal{O}_{F, \overline{z}}$ is a regular local ring
(Properties of Spaces, Lemma \ref{spaces-properties-lemma-regular}).
Note that a regular local ring is Cohen-Macaulay
(Algebra, Lemma \ref{algebra-lemma-regular-ring-CM}).
Let $d = \dim(\mathcal{O}_{F, \overline{z}})$. By
Algebra, Lemma \ref{algebra-lemma-find-sequence-image-regular}
we can find $f_1, \ldots, f_d \in \mathcal{O}_{U, u}$ whose images
$\varphi(f_1), \ldots, \varphi(f_d)$ form a regular sequence
in $\mathcal{O}_{F, \overline{z}}$. By
Lemma \ref{lemma-slice}
after shrinking $U$ we may assume that
$Z = V(f_1, \ldots, f_d) \to \mathcal{X}$ is flat and
locally of finite presentation. Note that by construction
$F_Z = Z \times_\mathcal{X} u$ is a closed subspace of
$F = U \times_\mathcal{X} u$, that $z$ is a point of this closed subspace,
and that
$$
\dim(\mathcal{O}_{F_Z, \overline{z}}) = 0.
$$
By
Morphisms of Spaces,
Lemma \ref{spaces-morphisms-lemma-dimension-fibre-at-a-point}
it follows that $\dim_z(F_Z) = 0$ because the transcendence degree
of $z$ relative to $u$ is zero. Hence it follows from
Lemma \ref{lemma-quasi-finite-at-point}
that after possibly shrinking $U$ the morphism $Z \to \mathcal{X}$
is locally quasi-finite.

\medskip\noindent
We conclude that for every finite type point $x$ of $\mathcal{X}$ there
exists a locally quasi-finite, flat, locally finitely presented
morphism $f_x : Z_x \to \mathcal{X}$ with $x$ in the image of $|f_x|$.
Set $W = \coprod_x Z_x$ and $f = \coprod f_x$. Then $f$ is flat, locally
of finite presentation, and locally quasi-finite. In particular the
image of $|f|$ is open, see
Properties of Stacks, Lemma \ref{stacks-properties-lemma-topology-points}.
By construction the image contains all finite type points of $\mathcal{X}$,
hence $f$ is surjective by
Lemma \ref{lemma-enough-finite-type-points} (and
Properties of Stacks, Lemma
\ref{stacks-properties-lemma-characterize-surjective}).
\end{proof}

\end{document}
